\chapter{Bevezetés}
% Csak példaként, hogy hogy működik a dolog:
\section{Szakmai háttér}
\paragraph{} Az elvégzett munka kontextusa, szakmai indoklása, motivációk, célkitűzés és az elvégzett feladatok összefoglalása.
\lipsum[5-6] % 2 random bekezdés - töröld ki

% -------------------------------------------------------------
\section{A szakdolgozat felépítése}
A szakdolgozat/diplomamunka felépítésének rövid összefoglalása:

\paragraph{2. Fejezet} Valamiről szól és most itt részletezem egy-két mondatban, hogy miről is van szó ebben a fejezetben. És most írok még egy mondatot.

\paragraph{3. Fejezet} Valamiről szól és most itt részletezem egy-két mondatban, hogy miről is van szó ebben a fejezetben. És most írok még egy mondatot.

\paragraph{4. Fejezet} Valamiről szól és most itt részletezem egy-két mondatban, hogy miről is van szó ebben a fejezetben. És most írok még egy mondatot.

\paragraph{5. Fejezet} Valamiről szól és most itt részletezem egy-két mondatban, hogy miről is van szó ebben a fejezetben. És most írok még egy mondatot.

\clearpage % Ez azért kell, hogy nehogy képek átcsússzanak a következő fejezethez